\hspace{3mm}

\centerline{\bf \Huge Домашнее задание.}
\centerline{\bf \Huge Арифметические прогрессии}

\begin{flushright}
    \scriptsize{Ты думаешь, что ты понял это, \\хотя на самом деле всего лишь думаешь,\\ что ты понял то, о чем ты думаешь. \\Хатаке Какаши. Наруто}
\end{flushright}

\problem{1!} Десятый член арифметической прогрессии равен 136, а двадцатый равен 536. Чему равен сотый член арифметической прогрессии?

\problem{2!} Пятидесятый член арифметической прогрессии на 2024 меньше сорок второго члена этой же прогрессии. Одиннадцатый член данной прогрессии равен нулю. Чему равен семьсот семьдесят седьмой член прогрессии?

\problem{3} За 17 февраля Бата сделал 100 прогибов, а за 21 февраля - 124 прогиба. За 18 февраля Фёдор спел 100 песен, а за 24 февраля - 148 песен. Каждый последующий день Бата делал на одно и то же количество прогибов больше, чем в предыдущий, а Фёдор пел на одно и то же количество песен больше, чем в предыдущий. Определите, кто из мальчиков был активней с 15 февраля по 29 февраля и на сколько, если 2 прогиба = 3 песням = 1 активности?

\problem{4} Известно, что сумма первого и девятого членов арифметической прогрессии меньше на 4, чем удвоенная сумма второго и четвертого членов этой же прогрессии. Сотый же член равен 1000. Найдите сумму первых пятидесяти членов данной арифметической прогрессии.

4